\section{Introduction}\label{sec:introduction}

An online \(k\)-server algorithm maintains \(k\) labeled servers in a metric
space. Upon receiving a request, it moves one server to the requested point
and leaves the others fixed. The movement is charged according to the metric.
The offline optimum knows the entire request sequence, whereas an online
algorithm must act before seeing the next request. For an initial labeled
tuple \(s\) and a finite request word \(\sigma\), write
\(\cost_{A,s}(\sigma)\) for an algorithm's total movement and
\(\OPT_s(\sigma)\) for the minimum such cost with full knowledge of the word.
Requests are fixed independently of the algorithm's random bits, the
oblivious-adversary model. We study finite rational metrics and the bit
complexity of each online decision.

The companion theorem~\cite{SquaredLog2026} supplies a policy with
competitive ratio \(O(\log^2(k+1))\), but no procedure for computing
its decision probabilities. We ask whether this existence guarantee can
be realized by a single uniform algorithm on all finite rational metrics,
using finite computations and unbiased random bits.

\begin{theorem}\label{thm:main}
There are an absolute constant \(C\) and a polynomial \(p\) with the following
property. Let \(n\ge3\), \(2\le k<n\), let \(d\) be a rational metric on
\(X=\{1,\ldots,n\}\), and let \(s\in X^k\) be an initial labeled tuple,
with repetitions allowed. Write \(L\) for the total binary encoding length
of \(n,k,d,s\). A single uniform randomized online algorithm preprocesses
these data in at most \(p(L)\) bit operations and serves request \(r_t\),
for every reachable history and internal state, in at most
\[
 p\!\left(L+\left\lceil\log_2(t+1)\right\rceil\right)
\]
additional bit operations. It uses only unbiased random bits and moves
exactly one chosen server to each request, permitting a zero-length move.
For every fixed finite request sequence \(\sigma\) independent of those bits,
\begin{equation}\label{eq:main-bound}
 \E\cost_{A,s}(\sigma)
 \le C\bigl(\log(k+1)\bigr)^2\OPT_s(\sigma)+B_{d,k,s},
\end{equation}
where \(B_{d,k,s}\) is finite and independent of \(\sigma\) and its length.
The algorithm does not need to know that length in advance.
\end{theorem}

Here and below \(\log\) is natural, except when the base is displayed.
The worst-case bit bounds in Theorem~\ref{thm:main} imply the corresponding
conditional expected bounds. The polynomial is uniform in every input
parameter and in \(t\).

Two features of the statement are essential. First, no bound on the size of
\(B_{d,k,s}\) is imposed. Second, the permitted computation grows polynomially
with the length of the binary request counter. Our algorithm serves
requests immediately, but may serve an enormous initial number of them
using one fixed label before the constructed policy takes over. The cost
of this initial service enters \(B_{d,k,s}\). Thus polynomial preprocessing
does not imply a short delay before activation or a polynomial bound on
the additive movement loss.

\subsection{History and computational scope}

Manasse, McGeoch, and Sleator introduced the \(k\)-server problem as a metric
framework for online service problems \cite{MMS1990}. Paging is its uniform
metric case: servers represent the \(k\) pages held in a cache, and a request
to an absent page costs one move. Fiat, Karp, Luby, McGeoch, Sleator, and
Young gave the randomized marking algorithm and the harmonic randomized
lower bound for paging \cite{FKLMSY1991}. This logarithmic behavior suggested
the classical randomized \(k\)-server conjecture, which proposed an
\(O(\log k)\) ratio on every metric. Bubeck, Coester, and Rabani disproved
that conjecture by establishing a worst-case lower bound of order
\(\log^2 k\) \cite{BCR2023}. The squared-logarithmic scale is therefore
necessary for a guarantee uniform over metrics.

Our mathematical input is Theorem~1.1 of the companion article
\emph{Squared-logarithmic randomized \(k\)-server on arbitrary metrics}
\cite{SquaredLog2026}. Its distinct-start clause supplies a single randomized
policy for all finite oblivious inputs with zero additive loss. We state
this exact finite-metric consequence as Theorem~\ref{thm:existence}.
Only the existence of that policy is used, not its probabilities or any
procedure for evaluating them.

Coester and Cosson address polynomial-time implementation directly
\cite[Theorem~1 and Corollary~2]{CoesterCosson2026}. They obtain ratio
\(O(\log^2 k)\) on sufficiently separated hierarchically separated trees
and ratio \(O(\log n\log^2 k)\) on arbitrary \(n\)-point metrics, with
polynomial request-processing time independent of the request-history
length. On the trees, they use only \(O(\log k)\) random bits in total;
the general-metric algorithm also uses a bounded initial bit budget
independent of the request-sequence length. These are stronger
time and randomness guarantees than those sought here. Theorem~\ref{thm:main}
instead obtains a competitive ratio independent of \(n\), using the
companion existence theorem and allowing request-processing time polynomial
in \(L+\log(t+1)\), together with an unrestricted finite additive movement
term.

Finite-policy synthesis and bounded-sequence reductions already have a
substantial history. M\"omke \cite[Sections~3--4]{Moemke2013} gives a
competitive-ratio approximation scheme on fixed finite metrics, reducing
the analysis to bounded request sequences and expressing randomized
strategies by a finite linear program. We use the same relationship between
causal probabilities and linear constraints, spelling out exact
elimination and the conversion to finite blocks of unbiased random bits.
Mikkelsen~\cite[Section~5.2, proof of Lemma~9]{Mikkelsen2016}, using
Lemma~8 there, develops related
reset constructions for finite lazy task systems: candidate policies are
simulated together, universally free requests are discarded, and a cost cap
triggers a fallback within a phase. Our universal filter and fallback
follow this general pattern, with explicit bounds for the retained word
and the finite state used by the implementation.
The proof retains the full construction so that its halting, sampling,
and storage properties can all be checked in the same bit model.

Resetting an online policy also requires accounting for the changes in its
server positions and for the cost of each completed phase. Such accounting
appears in M\"omke's bounded-sequence reduction
\cite[Section~3]{Moemke2013} and in the bounded-cost phase methods of
Komm, Kr\'alovi\v c, Kr\'alovi\v c, and M\"omke
\cite[Section~4]{KKKM2022}. Our fallback uses the protection of marked
points from classical paging \cite[Section~3]{FKLMSY1991}.

\subsection{Proof ideas}

The first step replaces real-valued policy probabilities by finite rational
data. For a fixed horizon, configuration probabilities and transition flows
form a rational linear system. Exact elimination finds a best feasible
competitive coefficient without being given the constant in the existence
theorem. Rounding the resulting conditional probabilities to dyadic values
adds a bounded expected movement error.

A finite horizon by itself does not control a long input whose optimum barely
moves. A service is \emph{strongly lazy} if it moves only on a request to an
unoccupied point. We retain a request only if some such service from the
fixed distinct starting tuple, using at most a prescribed number of positive
moves, would miss it. The filter represents these possible services by a
tree, branching into at most \(k\) choices at each miss and with depth
bounded by the move cap. Each retained request removes at least one live
node; no removed node returns. The resulting word therefore has bounded
length even when the original input does not.

A table trajectory that exceeds the move cap has already paid enough to
cover a deterministic marking fallback. An epoch groups a fixed number of
blocks, each completed block ending at its \((k+1)\)st distinct request.
Every completed block forces movement by the optimum, which pays for
restarting the table between epochs. Both the blocks and the filtered
request word are determined by the raw input alone. Consequently the
finite-policy estimate is always applied to a fixed word, even when the
actual service switches modes at a random time. This separates the finite
table's horizon from the potentially unbounded duration of an epoch.

Finally, a deterministic bit-machine constructor is simulated from scratch for only
\(\lfloor\log_2(t+1)\rfloor\) steps at request \(t\). Once it finishes, its
explicit output is small enough relative to that same budget to support all
subsequent online operations. Its finite completion time affects the additive
movement constant. Here the bounded retained word and explicit output-size
bound make that schedule possible despite the constructor's uncontrolled
running time.
This delayed activation is an implementation device made possible by the
unrestricted additive movement term; the substantive input to it is the
bounded-state suffix routine.

Section~\ref{sec:finite-tables} constructs the tables.
Section~\ref{sec:bounded-suffix} turns one such table into a routine for an
unbounded suffix. Section~\ref{sec:uniform-computation} schedules its
construction and proves Theorem~\ref{thm:main}.
